\documentclass{article}
\usepackage{graphicx}
\usepackage{amsmath}
\usepackage{hyperref}

\title{Room Acoustics: A Working Bibliography}
\author{Example Research Group}
\date{2026}

\begin{document}
\maketitle

\section{Introduction}
Modeling room acoustics spans geometric, wave-based, and data-driven methods
\cite{gauss2026,muller2025,zhao2025}. Perceptual studies motivate the choice of
metrics and rendering strategy \cite{alvarez2024,okafor2016}.

\section{Methods}
Our pipeline couples a sparse geometric model \cite{zhao2025} with a
physics-informed renderer \cite{gauss2026,gauss2026}. Reverberation time is
estimated directly from speech \cite{nasser2023}, and material parameters are
inferred with a diffusion model \cite{laurent2024}. See the measured impulse
response corpus for evaluation data \cite{park2022}.

\section{Related work}
Wave-based simulation at interactive rates is discussed by Müller and Example
\cite{muller2025}. Local wavefield synthesis provides an alternative spatial
rendering approach \cite{okafor2023}. For broader treatments, see the monographs
by Ivanov \cite{ivanov2022} and the auralization thesis by Example
\cite{example2021}. Practical geometric-acoustics guidance is collected in
\cite{muller2020}, while machine-learning approaches to acoustic parameters are
surveyed in \cite{chen2020}.

\bibliographystyle{plain}
\bibliography{library}

\end{document}
